El presente trabajo permitió especificar, implementar y validar exitosamente un modelo de simulación de eventos discretos (DEVS) para el controlador de una bomba de infusión intravenosa en lazo cerrado. El diseño demostró ser robusto frente a la naturaleza estocástica del entorno físico y operativo, mitigando los riesgos inherentes a un Sistema de Tiempo Real en el dominio médico.

En respuesta a los objetivos fundamentales del proyecto, el análisis automatizado de las trazas de simulación nos permite concluir que el modelo satisface con precisión matemática todos los requisitos y restricciones temporales especificados:

\begin{itemize}
    \item \textbf{Latencia de inicio y actuación:} Se comprobó empíricamente que toda nueva prescripción médica con caudal válido genera una respuesta física en el actuador en un tiempo estrictamente inferior a los 3 segundos. Asimismo, el mecanismo absorbe correctamente los retardos mecánicos uniformes del motor (hasta $0.5~s$).
    
    \item \textbf{Monitoreo y tolerancia a desvíos:} La arquitectura basada en un integrador temporal demostró ser altamente efectiva para evitar la fatiga por alarmas. El sistema toleró correctamente fluctuaciones gaussianas transitorias, emitiendo la advertencia de nivel medio (\textit{Alarma Media}) de manera exacta al acumularse 5 segundos continuos de un desvío superior al $10\%$.
    
    \item \textbf{Plazos de seguridad y detención crítica:} Ante fallas persistentes que alcanzaron el umbral de los 10 segundos sostenidos, el controlador logró ejecutar el bloqueo automático de la infusión, priorizando la integridad del paciente.
    
    \item \textbf{Ventana de gracia por fin de bolsa:} Las simulaciones confirmaron que el sistema detecta de forma temprana el vaciado de fluidos, emite la alerta preventiva y detiene automáticamente la bomba al cumplirse estrictamente el período máximo permitido de 60 segundos.
    
    \item \textbf{Vivacidad en la re-notificación:} Ante un escenario crítico de negligencia o ausencia del personal médico, el módulo de gestión demostró cumplir con los plazos de escalabilidad: aguardando inicialmente los 30 segundos requeridos tras una Alarma Crítica y entrando posteriormente en un bucle teóricamente infinito de repetición estricta cada 10 segundos.
\end{itemize}

En síntesis, la conjunción de las pruebas unitarias aisladas, los escenarios controlados (\textit{monkey-patching}) y la evaluación estocástica prolongada permiten garantizar formalmente que la lógica de control desarrollada opera siempre dentro de los plazos críticos de seguridad dictados por las normativas de la cátedra. Asimismo, la simulación del sistema y el análisis de los resultados obtenidos permitieron señalar posibles mejoras al diseño del sistema propuesto.
